\section{Où passe l'argent}
\label{sec:couts}

Le résultat affiché par le moteur n'est pas celui que le marché lui a donné~:
entre les deux, il y a des coûts. Cette section les mesure sur le journal des
transactions, poste par poste, et vérifie que la somme retombe exactement sur le
chiffre publié.

\begin{figure}[H]
\centering
\includegraphics[width=\linewidth]{gold_trades_attribution}
\caption{Du résultat de prix au résultat net. Les trois postes sont mesurés sur
le journal des transactions et se recollent exactement au chiffre publié.}
\end{figure}

% Généré par scripts/build_gold_trades_figures.py — NE PAS ÉDITER À LA MAIN.
\begin{table}[H]
\centering
\begin{tabular}{@{}lrl@{}}
\toprule
Poste & USD & Source \\
\midrule
Résultat de prix (mesuré) & +44,926 & CSV des deals \\
Swap (mesuré) & -12,784 & 28.5\,\% du brut \\
Commission (mesurée) & +0 & nulle sur ce compte \\
\textbf{Résultat net} & \textbf{+32,141} & recolle au rapport MT5 \\
\bottomrule
\addlinespace[0.4em]
\multicolumn{3}{@{}p{0.92\linewidth}@{}}{\footnotesize Slippage estimé à 545\,\$ (2\,bps par côté sur 1.4\,M\$ de notionnel cumulé) : il est déjà contenu dans les prix d'exécution, ce n'est pas un poste additionnel. Le stop de sécurité a coupé 1 position pour -1,955\,\$. L'arrondi de lots au pas de 0,01 n'est pas chiffrable depuis les artefacts disponibles et reste \emph{non attribué}.}
\end{tabular}
\caption{Décomposition du résultat : ce qui est mesuré, ce qui est estimé}
\end{table}


Le résultat de prix --- ce que le mouvement de l'or a rapporté --- s'élève à
$+$44\,926\,\$. Le portage en retire $-$12\,784\,\$, soit \metric{28,5\,\%}. La
commission est nulle sur ce compte. Il reste $+$32\,141\,\$, qui est exactement
le chiffre du rapport d'exécution~: l'attribution est complète, sans résidu.

\subsection{Le portage est un loyer de temps}

\begin{figure}[H]
\centering
\includegraphics[width=\linewidth]{gold_trades_costs}
\caption{Résultat de prix et résultat net cumulés~: l'écart entre les deux
courbes est le portage. À droite, le portage payé par trade en fonction de sa
durée.}
\end{figure}

Le panneau de droite explicite le mécanisme~: le portage croît avec la durée de
détention, et frappe donc le plus lourdement les positions qui rapportent le
plus. Le trade de 214~jours paie 2\,874\,\$ de portage, soit 11,7\,\% de son
gain de prix~; celui de 199~jours en paie 2\,844\,\$, soit 18,9\,\% du sien.

Rapporté à l'ensemble, le portage représente pourtant 28,5\,\% du résultat de
prix --- nettement plus que sur les grands trades. La raison est simple~: les
positions perdantes paient elles aussi leur loyer, sans rien rapporter en
face.

\begin{insightbox}
C'est le coût structurel de la méthode, et il ne peut pas être supprimé sans
supprimer ce qui fait gagner le moteur. Raccourcir les positions réduirait le
portage, mais amputerait d'abord les gains~: rappelons que les cinq positions
de plus de cent jours produisent \emph{plus} que le résultat total.
\end{insightbox}

\subsection{Deux postes plus fins}

Deux coûts complètent le tableau, et il faut être clair sur leur statut.

Le \textbf{glissement} --- l'écart entre le prix visé et le prix obtenu --- est
estimé à environ 545\,\$ sur 1,4~million de dollars de notionnel cumulé, à
partir du paramètre de 2~points de base par côté. Cette somme est déjà contenue
dans les prix d'exécution du tableau ci-dessus~: ce n'est pas un poste à
retrancher en plus.

L'\textbf{arrondi des volumes} au pas de 0,01~lot n'est pas reconstructible à
partir des artefacts disponibles. Il est réel, son ordre de grandeur est discuté
au chapitre~\ref{sec:sizing}, mais il reste \emph{non attribué} --- et ce
document ne lui invente pas de chiffre.
